\begin{tikzpicture}[
    compedge/.style={edge, draw=red}
]

% Oba grafy maji vrcholy ve stejnych pozicich: pravidelny osmiuhelnik,
% vrcholy 1 az 8 po smeru hodinovych rucicek, zacina se vlevo nahore.
% Zadne tri vrcholy nelezi na jedne primce, takze je videt kazda hrana.

%================================================
% G = Wagneruv graf: kruznice C_8 a ctyri spojnice protilehlych vrcholu
% (kazdy vrchol ma stupen 3)
%================================================
\begin{scope}[xshift=0cm, yshift=0cm]

    \foreach \i/\uhel in {1/112.5, 2/67.5, 3/22.5, 4/337.5,
                          5/292.5, 6/247.5, 7/202.5, 8/157.5} {
        \coordinate (a\i) at (\uhel:1.7);
    }

    % kruznice C_8
    \foreach \u/\v in {1/2, 2/3, 3/4, 4/5, 5/6, 6/7, 7/8, 8/1} {
        \draw[edge] (a\u)--(a\v);
    }

    % spojnice protilehlych vrcholu
    \foreach \u/\v in {1/5, 2/6, 3/7, 4/8} {
        \draw[edge] (a\u)--(a\v);
    }

    % vrcholy
    \foreach \i in {1,...,8} {
        \node[vertex, fill=white] at (a\i) {};
    }

    \node at (0,-2.3) {$G$};

\end{scope}

%================================================
% Doplnek \bar{G}: vsech 28 - 12 = 16 zbyvajicich hran
% (kazdy vrchol ma stupen 7 - 3 = 4)
%================================================
\begin{scope}[xshift=5.2cm, yshift=0cm]

    \foreach \i/\uhel in {1/112.5, 2/67.5, 3/22.5, 4/337.5,
                          5/292.5, 6/247.5, 7/202.5, 8/157.5} {
        \coordinate (b\i) at (\uhel:1.7);
    }

    % vrcholy ob jeden (dva ctverce)
    \foreach \u/\v in {1/3, 2/4, 3/5, 4/6, 5/7, 6/8, 7/1, 8/2} {
        \draw[compedge] (b\u)--(b\v);
    }

    % vrcholy ob dva (osmicipa hvezda)
    \foreach \u/\v in {1/4, 2/5, 3/6, 4/7, 5/8, 6/1, 7/2, 8/3} {
        \draw[compedge] (b\u)--(b\v);
    }

    % vrcholy
    \foreach \i in {1,...,8} {
        \node[vertex, fill=white] at (b\i) {};
    }

    \node at (0,-2.3) {$\bar{G}$};

\end{scope}

\end{tikzpicture}
